\chapter{가설 충족 현황과 한계}

\section{가설별 현황}

다섯 가설 가운데 H1·H4는 해당 모델 범위에서 충족했고, H3는 명시한 기동·부하 조건(사전 기동, 600 µs 램프, 보조부하 25 kW)에서 충족했으며(640 V·0·1 kW 기동은 불통과), H2는 부분 충족(640 V 복합 공차 Tier-2 평균모델 102/128), 기준 구조 대비 이점(H5)은 기각되었다. H5 기각 뒤 Q3에 더한 두 기준(적용 범위, 산업 방식 대비)도 함께 적는다(표 \ref{t:status}). 모든 판정은 해당 모델과 조건 안에서만 유효하며, 하드웨어 성능을 뜻하지 않는다.

{\footnotesize
\begin{longtable}{L{\dimexpr 0.13\linewidth-2\tabcolsep\relax}L{\dimexpr 0.19\linewidth-2\tabcolsep\relax}L{\dimexpr 0.15\linewidth-2\tabcolsep\relax}L{\dimexpr 0.38\linewidth-2\tabcolsep\relax}L{\dimexpr 0.15\linewidth-2\tabcolsep\relax}}
\caption{가설과 연구 질문 Q3의 최종 판정}\label{t:status}\\
\toprule
\tabhead \textbf{가설} & \textbf{수락 기준} & \textbf{판정} & \textbf{근거} & \textbf{중간보고서 이후 바뀐 점} \\
\midrule
\endfirsthead
\toprule
\tabhead \textbf{가설} & \textbf{수락 기준} & \textbf{판정} & \textbf{근거} & \textbf{중간보고서 이후 바뀐 점} \\
\midrule
\endhead
\bottomrule
\endlastfoot
H1 정격 분리 & 링크 ≤ 8 kW(정상·Tier-1·Tier-2, 공차 포함), 충전전류 Stage B 비경유 & \textbf{충족(Stage B가 있는 구조)} & 주행 정상 6.44 kW, 주행 과도 정식 7.16 kW(E5 640–920 V 516건; E7 복합 공차 256건 7.12 kW), 충전 중(640–920 V) 정적 6.29 kW·과도 정식 7.04 kW(516건; 632 V 경계 진단 포함 6.69·7.49 kW). 충전전류는 V\textsubscript{A} 노드에서 Stage A로(E3, E5–E7, E16, TB07, E36) & A25R에서 다시 확인 \\
H2 강건 제어 & Tier-1·Tier-2, ±10 \% 공차, 추가 지연 +3.0 µs & \textbf{부분 충족(Stage B 모델 범위)} & Tier-1은 공차·지연 포함 평균모델·독립 구현·스위칭 모델 모두 통과. Tier-2는 단일 공차 격자 78/78, 640 V 복합 공차 128개 조합에서 평균모델(93.75 A 기준) 102/128(최대 95.82 A)·스위칭 모델(순시 100 A 기준) 128/128(97.32 A). 단일 공차 격자의 +3 µs 정보 시험 13/14(632 V 경계 94.08 A; A25R E36도 13/14). 창 조정 모듈 456/456·264/264(E5–E7, E11, E12, TB04–TB08) & 복합 공차 Tier-2 평균모델 26/128 불통과 \\
H2 추가 시험(E46) & 같은 128개 조합의 스위칭 모델(순시 100 A 기준)에 H2 기준의 추가 지연 +3 µs & \textbf{불통과 14/128(H2 부분 충족에 포함)} & 114/128(실패 순시 100.05–101.21 A, 링크 최대 6.83 kW; 같은 128개 조합에 추가 지연 +3 µs를 준 별도 시험이며, 지연 미추가 128/128과 합산하지 않음) & 두 번째 통합 검토에서 추가 \\
H3 충전 중 보조버스 & 150 kW 충전 + 25 kW 보조부하 동시, 접속·해제 전이(충전기 명령 I = min[375 A, 150 kW/(V\textsubscript{bat}/2)]) & \textbf{충족(조건: 충전 전 사전 기동·충전전류 램프 ≥ 600 µs, 보조부하 25 kW)} & E16, TB07, E26, TB20–TB22, E35, E36, TB27, TB28. 600 µs는 보조부하 15·25 kW에서 확인했고, 640 V·0–1 kW는 위상 첨두 78.63–78.74 A로 E35 공칭 한도 74.53 A 초과(E46, 별도 기록) & 조건 100 → 600 µs, 사전 기동 추가 \\
H4 부하 승인 & 예산 초과 0, 승인 지연(무손실 최대 ≤ 25 ms, 손실 ≤ 20 \%에서 99백분위 ≤ 50 ms), 미공급 < 정적 예약 & \textbf{충족(축약 모델)} & 예산 초과 0 ms, 승인 지연 무손실 최대 17 ms, 손실 조건 99백분위 25.44 ms·최대 32 ms, 미공급 6.96 \%(정적 예약 35.21 \%, E8). 승인 지연 모집단은 200개 × 20 s 축약 시나리오의 압축기 승인 488건(무손실 235·손실 253)이며, 연속 통신 단절·버스 혼잡·모든 차량 부하의 확률 보장은 주장하지 않는다(E46). 고장 강등 연동 74/74(E29) & 보호 상태기계 연동 추가 \\
H5 기준 구조 대비 이점 & 손실·정격 VA·반도체 수 중 2개 이상 우수, 서비스 동등 이상 & \textbf{기각} & ③ 144.5 대 ④ 101.5 W, 52 대 42위치, 1,099 대 1,009 kVA(E15, E36, E38) & A25R·최종 Stage C로 다시 계산해도 같음 \\
Q3 적용 범위(추가) & 성립 구조 가운데 최소 손실인 팩·부품 조합 & \textbf{창 조정 모듈은 해당 없음(같은 수준)} & 공정 기준에서 창 조정 모듈이 최선이던 세 칸 모두 다른 성립 구조와 같은 수준(E42). ①은 B\textsubscript{420} 216·224s·사양 팩에서 최선 & 이전 판정(창 조정 모듈 최선)에서 바뀜 \\
Q3 산업 방식 대비(추가) & R2가 앞서는 조건과 뒤지는 조건 & \textbf{제한적} & 충전 경로 손실만 이점(전용 승압기의 0.32–0.51배), 보조 손실은 ⑤와 같은 수준·⑥보다 4–33 W 큼(E41, E43, E44) & 이전 판정(⑤보다 12–26 \% 작음)에서 바뀜 \\
\end{longtable}
}
\srcnote{주: H2의 회로 확인은 Stage B 모델 범위다. H3의 최종 설계(A25R) 근거는 Stage A 정적 등가로 다시 계산한 사례 집합(E36)과 결합 모델·회로의 대표 사례(E35, TB27, TB28)이며, 정적 등가의 보수성은 증명하지 않았다(18.4절). 충전기별 실제 램프 응답은 하드웨어 단계의 확인 항목이다(19장).}

\section{불통과와 부정적 결과}

판정 기준은 단계마다 시험 전에 정하는 것을 원칙으로 했다. E30 이후의 시험은 기준과 코드를 실행 전에 기록해 두었으나, E26–E29와 TB21–TB23은 기준과 결과가 같은 기록에 있어 순서를 기록으로 확인할 수 없다(4장). 설계를 고르는 데 쓴 사례가 판정 사례에도 들어간 경우(A25R의 느린 충전 3·Tier-1 계단 2, 13장)는 따로 표시했다. 불통과 결과는 지우지 않으며, “처리”는 그 결과로 바꾼 설계나 남긴 과제다.

\Needspace{8\baselineskip}
{\footnotesize
\begin{longtable}{L{\dimexpr 0.17\linewidth-2\tabcolsep\relax}L{\dimexpr 0.47\linewidth-2\tabcolsep\relax}L{\dimexpr 0.36\linewidth-2\tabcolsep\relax}}
\caption{불통과·부정적 결과와 처리}\label{t:fails}\\
\toprule
\tabhead \textbf{시험} & \textbf{결과} & \textbf{처리} \\
\midrule
\endfirsthead
\toprule
\tabhead \textbf{시험} & \textbf{결과} & \textbf{처리} \\
\midrule
\endhead
\bottomrule
\endlastfoot
E1 Stage A 고정 타이밍 & 20건 중 16건 영전류 스위칭 실패(잔류 최대 40.1 \%) & 영교차 추종 게이트 채택 \\
\textbf{E7 복합 공차} & 640 V 128개 끝값 조합의 Tier-2가 평균모델 93.75 A 기준 102/128(최대 95.82 A; L\textsubscript{h} ±5 \%에서도 8/26 남음), 스위칭 모델 100 A 기준 128/128 & 기준 유지, H2 부분 충족 \\
E9 제어기 대안 & 대역폭 저하·전압항 포화로는 지연 여유 목표 미달 & 예측 보상 채택 \\
E13 비조정 구조 ④ & 경로 L–버스 C 공진으로 Stage A 전류 최대 136.8 A(Tier-1 462/490, Tier-2 56/98) & 버스 RC 감쇠기(1.6 mF·48.6 mΩ) \\
E14 Stage B 정지 & 직렬전압이 남은 채 단락하면 571–1,083 A & 보호·기동 요구 \\
\textbf{E15 H5} & \textbf{기각}: ③이 ④보다 손실·반도체 수·VA 모두 큼 & Stage B를 선택 모듈로, R2로 구조 변경 \\
E16·TB07 & 충전전류 10 µs 변화에서 경로 첨두가 등가 모델보다 3.3–13.4 A 큼(최대 98.49 A) & 충전 램프 요구(최종 600 µs) \\
E19 & 창 조정 DAB 정격 1.25배에서 포화, 즉시 바이패스 740–869 A & 정격 1.8배, 강등 순서 \\
TB09 1차 & H-브리지 개방 2사례 불통과(링크 정체, 소프트 투입 177 A), DAB 순환전류 정격점 4.2배 & 정체 감지·경로전류 조절, EPS 변조 \\
E17 변조율 고착 & 0/16(필터 128–147 A) & 필터 인덕터 포화 전류 ≥ 150 A 요구 \\
E20 정책 비교 & 기존 추천 0/8, EPS만 0/4, 감발 없음 3/4 & 정합 유지 + EPS 첨두 상한 + 일시 감발 \\
TB10·\textbf{TB11}·\textbf{TB12} & 1.5배 목표 초과(붕괴 구간 10–16 \%), 지령 전이 직류 편류 88.5 A, 전이 끝 1.3 A 초과(1/2) & 데드타임 반영 첨두 표, 편류 방지, 설계 상한 1.387배(TB13 2/2) \\
\textbf{TB14} & 1차 0/5(영전압 스위칭 실패), 2차 6/7(역방향 경부하 이득 상승) & 레그 출력용량 모델로 재설계, 경부하 두 브리지 구동(TB15 5/5). 역방향 0.6–0.74 kW 이득 0.96–0.97은 남음 \\
E23 Stage A 경부하 & 1008 V 강등 뒤 10 A에서 턴오프 비율 5.40–5.47 \% & 위상 셰딩 \\
E24 사례 E & 충전 램프 요구의 4배에서 위상 첨두 93.6 A & 요구 밖으로 기록 \\
\textbf{TB20 사례 D} & 턴오프 1번 1.010 A(기준 1.0 A) & 게이트 방향 반전 처리(TB21 4/4) \\
E27 회귀 & 엄격 기준 불통과(판정 변화 없음, 지표 차 두 가지) & 그대로 기록 \\
\textbf{선행 대조} & “같은 2:1 단 계속 사용”은 개념의 선행이 있음(중간보고서와 불일치) & 기여 1 수정(15장) \\
E28 & 1·2차 실행 불통과, 정격 충전점 시작 과도 12점, 사전 기동 없는 충전 84.2 A(회로 83.2 A) & 설계 변경 4가지, 사전 기동 요구, 판정 창 그대로 \\
\textbf{E29} & 48 V 회생 차단 없음 0/4(버스 432.7–527.0 V), 이중 고장 경로전류 5 ms 1/4 & Stage C 회생 차단 요구, 이중 고장은 E38에서 설계 변경 \\
\textbf{E31} & \textbf{Stage A 출력용량 손실 항이 한 자릿수 작았다}(원설계 위상당 175 W, 모델 14 W). 원설계에서는 R2 손실 결론이 모두 뒤집힘 & 공진주파수를 낮춘 설계(25 kHz) \\
TB25 & 전체 9개 중 3개 불통과(무부하 3/3, 부하 대조 3/6: T2 2/3·T3 1/3) & 그대로 기록(무부하 손실은 식과 같음) \\
\textbf{E33·TB26} & A25 과도 1/7, 느린 충전 0/3, 회로 3/3 불통과 & C\textsubscript{A} 6배(A25R) \\
E35 & 설계 단계 후보 대부분 불통과, 요구 밖 조건 0/12 & A25R 선택, 요구 밖으로 기록 \\
\textbf{E34} & 부품 1차 9개 중 4개 불통과 & 2차·3차에서 설계로 해결 \\
\textbf{E36} & 920 V 충전이 모터 승압 하한보다 6 W 큼, B\textsubscript{420} 두 칸 바뀜, 이중 고장 0/4 & 문장 정정, 400 V 부하 정지 \\
\textbf{E37} & 소프트 투입 선형 MOSFET 245.5 °C 이상, 방전 저항 2/4 & E38 설계 변경 \\
E38 & 클램프 한 가지 저항 불성립, 5 ms 정지 램프에서 직렬 114.0 V, 240s R2-W460 VA가 ②보다 큼 & 가지 분할, 3 ms 선정, 문장 정정 \\
\textbf{E39} & E38 클램프 스위치 183.3·180.8 °C 불통과(전형 열저항 구현 누락) & CAB530M12BM3 \\
\textbf{E40} & E39 클램프 159.7·157.6 °C 불통과(규칙 L), 최대 켜짐 시간 없으면 오동작 클램프가 100 ms 켜져 11.8–12.6 kJ, 문턱 +0.15 \% 이상 C1 불통과 & 가지 3개 × 2.4·2.7 Ω + FF2MR12KM1P, 최대 연속 켜짐 0.5 ms, 문턱 정확도 요구 \\
\textbf{E41} & ④ 대 ⑤가 공정 기준 0.5 \%, 민감도로 부호 바뀜 & 결론 2·기여 1을 “같은 수준”으로 \\
\textbf{E42} & ①의 보조 강압기에 셰딩을 준 공정 기준에서 창 조정 모듈이 최선이던 칸(B\textsubscript{420} 200s, B\textsubscript{470} 224·240s)과 ④가 최선이던 칸이 모두 “같은 수준” & 기여 2·Q3에서 창 조정 모듈의 보조 손실 이점 주장 철회, 성립 수단으로 위치 \\
\textbf{TB33 C33} & 오동작 사례 평균모델 대조 불통과(회로가 9 µs 만에 끔). 진단에서 클램프가 길게 켜질수록 Stage A 출렁임(위상 첨두 155.2 A, 턴오프 초과 최대 20번) & 이 사건만 평균모델이 에너지를 크게 봄을 기록(일반적 보수성 근거 아님), 진단은 판정 아님 \\
\textbf{E43} & Stage C 일시 정지 동안 48 V는 배터리 전압(20 mΩ 등가 41.6–45.7 V)으로 ±2 \% 미유지, 60 mΩ(OCV 44 V)에서 35.6 V & 48 V 연속성 범위 축소, 배터리 요구(≤ 57.6 mΩ) 도출, 결합 시험은 후속 \\
\textbf{E44} & E\textsubscript{oss} 조건 정정으로 ④ 대 ⑤ 명목 순서 바뀜(⑤ 1.4 \% 작음). 클램프 +0.10 \%·수동소자 −10 \%에서 버스 470.08 V, 폭 3.0 V에서 150.6–153.9 °C, W460 3.5 V 모서리 150.7 °C & ‘같은 수준’ 유지, 구현 요구 +0.075 \%·폭 3.5(W460 4.0)–6.0 V \\
\textbf{E45} & 클램프 지연 0.5 µs·폭 3.5 V에서 151.26 °C(반례 9건). 가치 지도 판정 함수의 명목 최소 비교. Windows spawn의 설정 상실 경로 & 하드웨어 목표·검증 입력 구분과 후속 시험, 200s/B\textsubscript{470} ①·④로 정정, 코드 수정 \\
\textbf{E46} & +3 µs 복합 공차 스위칭 114/128. 640 V·0·1 kW 충전 시작 78.63·78.74 A(한도 74.53 A; C\textsubscript{fly} 후보 76.06·76.56 A). 216s 충전 단자를 이전 정의로 계산. 검증기 우회 & H2 추가 기록, H3 부하 조건, 직결 단자 150.166 kW, 검증기 수정 \\
\end{longtable}
}
\srcnote{주: 굵은 시험은 설계나 결론을 바꾼 결과다. 실행 경과(설정 오류, 구현 누락의 수정, 1차 결과 보관)는 각 근거 문서에 모두 적었다(부록 B).}

\Needspace{10\baselineskip}\section{검토에서 나온 지적과 대응}

{\footnotesize
\begin{longtable}{L{\dimexpr 0.20\linewidth-2\tabcolsep\relax}L{\dimexpr 0.44\linewidth-2\tabcolsep\relax}L{\dimexpr 0.36\linewidth-2\tabcolsep\relax}}
\caption{검토 지적과 대응}\label{t:reviews}\\
\toprule
\tabhead \textbf{검토} & \textbf{지적} & \textbf{대응} \\
\midrule
\endfirsthead
\toprule
\tabhead \textbf{검토} & \textbf{지적} & \textbf{대응} \\
\midrule
\endhead
\bottomrule
\endlastfoot
외부 검토(2026-10-29) & 6곳 재현으로 타당(손실 비교의 Stage C, 클램프 저항 공칭값·IGBT 손실 온도·설정 일괄 적용, “150 kW” 표기, 충전 손실 제곱 환산) & E38, TB30·TB31 \\
재검토(2026-10-29) & 클램프 스위치를 전형 열저항으로 판정, 가지 저항 섞임 미판정 & E39 재판정, TB32 \\
6인 검토(2026-10-29) & R\textsubscript{DS(on)} 등 전형값 의존, ④ 대 ⑤ 비교 조건 불공정, 클램프 오동작 실패 모드, C\textsubscript{A} 후보 표 잘림 & E40·E41·TB33 \\
최종 보고서 독립 검토(2026-11-03) & 가치 지도의 ①이 보조 강압기를 셰딩 없이 둠(비교 조건 불공정) & E42 재계산, 창 조정 모듈의 보조 손실 이점 주장 철회 \\
제출 전 외부 검토(2026-11-05) & 버스 측정 범위, 48 V 유지 근거, 정적 등가 보수성, 민감도 스위칭 배율, W460 자기 가열 가정, 문구·규격 & E43(7.4·10.4·12.6절, 표 \ref{t:e41}) \\
제출 전 재검토(2026-11-06) & E\textsubscript{oss} 적용 전압, H4 조건, H1 632 V 혼입, 클램프 수동소자·추종 오차, C\textsubscript{fly} 후보값 & E44(15.2절, 12.6절, 16장) \\
제출 전 통합 수정안(2026-11-06) & H2 실패 누락, 클램프 지연 반례, P5 온도 정의, 관찰 종료점, 노드별 상한·충전전력, 가치 지도 판정 함수·180s, 실행 방식, 스위치 집계 & E45(docs/50) \\
제출 전 두 번째 통합 검토(2026-11-06) & H2 +3 µs, H3 경부하 기동, 충전 단자, W410 1차 전류, P6 한도, TB25 분모, 정지 에너지 포트, CV 범위, 교차점 기준, 검증기 & E46(docs/51) \\
최종 제출 판정(2026-11-07) & 시험 종류 혼합 서술, 부하 분리 조건 문단, 충전 과도 분모·스텝, E11 표현, 그림 2(a) 축, 검증기·해시 & E47(docs/52), 원고 분리 서술 \\
\end{longtable}
}

\section{현재 근거의 한계}

\begin{itemize}
\item \textbf{하드웨어 시험이 없다.} 모든 판정은 해석·평균모델·스위칭 모델·소자 수준 회로 모델·데이터시트 대조 안에서만 유효하다. 회로 모델의 스위치 셀은 이상 스위치, 다이오드와 선형 출력용량이며, 게이트 드라이버, 배선 기생성분, 열, 절연과 전자기 간섭은 넣지 않았다.
\item \textbf{손실 비교는 같은 규칙의 상대 비교다.} 출력용량 항의 규칙은 구조마다 다르다(Stage A는 Q\textsubscript{oss}·V, 하드 스위칭 변환기는 E15 규칙). 측정된 2:1 하드웨어(100 kW, 97.4 \%)와 모델의 차가 크며(15.2절), TB27의 +3.0 \%는 정격 충전점 한 회로 사례(510.5 대 495.8 W)라 손실 모델 전반의 정확도로 인용하지 않는다. 평균모델의 DAB 고정 효율 가정은 구조 비교에 쓰지 않는다. 모터 승압 손실은 리플 철손·자석 손실이 없어 하한 쪽이고, ⑤ 25 kHz 경우는 ⑤에 유리한 끝값이다. 가치 지도(E42)는 Stage A 고정 손실 민감도를 넣지 않았다. 명목 손실은 정적 모드 선택의 대표 운전점 가중 평균이며(히스테리시스·셰딩 유지시간은 단일 경계로 근사), 실주행 시계열 손실은 계산하지 않았다. 비교 소자 E\textsubscript{oss} 적용 전압은 데이터시트 곡선과 대조해 다시 정했고(E44), GS66516T 스위칭 에너지는 25 °C 표 값, EPC2305·GS66516T 역도통 전압과 EPC2305의 125 °C 온저항 배율은 곡선에서 정한 가정값이다.
\item \textbf{A25R로 다시 확인하지 않은 것.} 멈춘 모듈 C\textsubscript{fly} 전압 요구(−10–0 V), 225 kW(12–16위상) 충전의 회로, Stage B 보호·강등 시험(E17, E19–E21; V\textsubscript{A}를 이상 전원으로 둠. 이중 고장과 소프트 투입만 다시 확인)이다.
\item \textbf{Stage A 정적 등가.} V\textsubscript{A} 노드 공진보다 느린 외란에서만 맞으며, C\textsubscript{A} 2.4 mF에서는 C\textsubscript{A}·R\textsubscript{A} 시정수(5–22 µs)가 제어주기보다 길다. E36의 정적 관계가 모든 판정량에서 보수적이라는 증명은 없고 TB07·TB33은 정적 등가에 없는 순간 동특성을 보였으므로, H3의 최종 설계 결론은 정적 등가의 사례 집합과 결합 모델·회로의 대표 사례(E35 과도 6/6(설계 선택 5사례 포함)·모듈 고장 16/16, TB27·TB28 4/4)로 한정한다.
\item \textbf{새 인터페이스 요구의 실제 응답.} Stage C 입력 제한(4 kW/ms 정지·재개, 제어기 미구현)과 400 V 부하의 3 ms 정지는 모델의 가정이다. 일시 정지 동안의 48 V는 배터리 등가로만 계산했고(E43, 10.4절), 배터리 인덕턴스·온도·SOC와 Stage C 재개 과도의 결합 시험은 하지 않았으므로 H-브리지 고장 강등에서 48 V ±2 \% 유지를 주장하지 않는다.
\item \textbf{고전압 팩(224·240s)의 충전 중 정지}는 V\textsubscript{A} 수준 모델만 확인했다(Stage B 결합·회로 없음). R2-W410 회로는 평균모델·스위칭 모의만 있다.
\item \textbf{회로 시험의 시간 축소와 이상화.} 400 V 부하 분리 20 → 3 ms(TB09–TB19), 충전기 정지 30 → 3 ms(TB20·TB21·TB28–TB30·TB32·TB33, 흡수 에너지는 E26), SALS 시간 0.15배(TB23·TB31)로 줄였고 기준값도 같은 축소로 계산했다. 회로의 충전기는 이상 전류원(충전기 커패시터·전압 루프·케이블 L 없음), 400 V 부하는 계약을 지키는 정전력 등가(스위칭 맥동 없음), 클램프 스위치는 이상 스위치, 부품값은 공칭이다. TB16–TB19의 Stage A는 모듈 하나(2위상)다.
\item \textbf{공차.} 공차 표본은 3개(A25R은 공칭·표본 1)이고 최악 조합 탐색·대규모 몬테카를로가 아니다. C\textsubscript{fly} 후보 45.0 µF와 −5 \%는 요구 안 과도·모듈 고장만 다시 판정했고(E44, 타이밍 45.0 µF 기준; 46.875 µF 타이밍과 45.0 µF 부품의 조합은 미실행), 위상 간 편차 최악 조합·부품 온도·요구 밖 과도는 넣지 않았다. 공차 표본에서 C\textsubscript{A} 실효전류는 인터리빙 상쇄가 깨져 12배가 되었다.
\item \textbf{Stage A 영전압 스위칭 보조 운전은 시험하지 않았다.} “150·100 kHz 0/30”은 조사한 후보 90개·영전류 게이트·당시 E31 동시 기준의 결과이며 고주파 운전의 불가능을 뜻하지 않는다(13장).
\item \textbf{부품 정격의 기본 판정은 데이터시트 표·본문의 수치로 했다}(정상 열저항 보수 계산). 손실 계산의 일부 소자 특성(비교 소자 E\textsubscript{oss}의 적용 전압, EPC2305·GS66516T의 역도통 전압과 온저항 배율)은 해당 항목에 밝힌 특성곡선 판독값이다. 규칙 L의 × 1.2는 정한 값이며, 실제 최대/전형 비가 더 크면 결과가 바뀔 수 있다(선정 변형의 여유 FF2MR12KM1P 24.9 °C, CAB530M12BM3 18.6 °C). 조건부 부품은 16장에 있다.
\item \textbf{V\textsubscript{A} 클램프.} 문턱 구현 요구(표 \ref{t:e26r})는 버스 수동소자 ±10 \%의 두 상관 모서리에서만 확인했고 여유가 작다(상한 쪽 0.037 V, 최소 폭 쪽 2.36 °C, 0.25 µs 시간 격자). 켬·끔 지연 1.0 µs 고정에서 K1·K2는 이산 탐색, K3는 지정 끝값의 모서리이므로 연속 범위 전체의 통과가 아니며, 0.5 µs에서는 사양 팩 3.5 V 폭이 151.26 °C로 넘었다(E45, 12.6절). 소자별 반대 방향 공차 조합, K2의 수동소자 공차, 검출 회로(분압기·절연 증폭기·비교기)의 설계·온도·노화 오차, 버스 배선 저항 공차, 잡음이 계속되는 반복 오동작, R2-W460 강제 끔 순간의 di/dt·스위치 과전압(이상 스위치 모델)은 보지 않았다.
\item \textbf{선행 대조}는 데이터베이스 전수 검색·청구항 대조가 아니다(15.3절).
\item \textbf{H2의 추가 지연은 복합 공차에서 불통과가 있다}(+3 µs 스위칭 114/128, E46, 5장). 센서·게이트 지연 범위를 정해 제어기를 다시 조정하는 일은 후속 과제다.
\item \textbf{충전 시작 조건은 보조부하에 따라 다르다.} 사전 기동·600 µs 램프는 보조부하 15·25 kW에서 확인했고, 640 V·0–1 kW에서는 위상 첨두가 공칭 한도를 넘는다(8.3절). 실차 충전 시작은 부하·온도·초기 모듈 수에 따른 기동·감발 행렬로 다시 정해야 한다.
\item \textbf{고장 강등 74/74의 범위.} 74개 가지는 지정한 차단 가능 고장과 메시지 지연·손실 조건이다(직접 분리 38, 감발 36 중 복원 25·분리 11). 임의 부품 단락·게이트 고착·통신 폭주는 조사하지 않았고, 48 V 회생 중 배터리의 충전 수용·BMS 동작은 결합하지 않았다.
\item \textbf{가치 지도의 이력 의존.} 정적 단일 모드 경계로 만든 표라, 200s/B\textsubscript{470}에서 W410의 후보 여부는 ±5 V 제어기의 이전 모드에 따라 달라질 수 있다(표 \ref{t:e18} 각주 b).
\item \textbf{표준은 설계 입력으로만 썼다.} ISO 21498-1:2021의 전압 하위 등급, ISO 21780:2020의 48 V 전압 범위, IEEE Std 2030.1.1-2015 부속서 A의 CHAdeMO 방식 시험 항목을 인용한 판본·방식·항목으로 한정해 요구로 옮겼다. 2015년 판은 IEEE에서 2021년 판으로 대체됐으며, 이 보고는 현행판과의 차이를 대조하지 않았다. 과도 중 420·470 V 상한은 표준 값이 아니라 이 과제의 설계 계약이다. 차량이나 부품의 인증·규격 적합성은 주장하지 않는다.
\end{itemize}
